% source: https://github.com/pfradin/luadraw/discussions/276#discussioncomment-17211790 (pfradin, block 1)
% https://tex.stackexchange.com/q/569439
% https://github.com/pfradin/luadraw/discussions/276#discussioncomment-17210452
\documentclass{standalone}
\usepackage[3d]{luadraw}
\usepackage[svgnames]{xcolor}
\usepackage{fourier-otf}

\begin{document}
\begin{luadraw}{name=sphere-cone}
local ld = luadraw
local pt3d,deg = ld.pt3d,ld.deg
local M, Ms = pt3d.M, pt3d.Ms
local O, vecI, vecJ, vecK = pt3d.Origin, pt3d.vecI, pt3d.vecJ, pt3d.vecK

local oldDsphere = ld.graph3d.Dsphere
function ld.graph3d:Dsphere(C,R, options)
    self:Filloptions("none","black")
    oldDsphere(self,C,R,options)
end

require 'luadraw_spherical'

local g = ld.graph3d:new{
    window = {-2,12.5,-2,7},
    size={12,12},
    viewdir={20,75}
}
local C,R,L = M(3,6,5), 3.1, 17 --M(1.8,7.2,5.4), 3.1, 17
local theta_start, theta_end = 60.0*deg, 71.0*deg
local phi_start, phi_end = 58.5*deg, 66.0*deg
local P1 = Ms(L,theta_start,phi_start)
local P2 = Ms(L,theta_end,phi_start)
local P3 = Ms(L,theta_start,phi_end)
local P4 = Ms(L,theta_end,phi_end)

local function hit(d)
    local A, B = ld.interDS({O,d},{C,R})
    if pt3d.dot(A-O,d) > pt3d.dot(B-O,d) then A, B = B, A end
    return A, B
end

local A1, B1 = hit(P1)
local A2, B2 = hit(P2)
local A3, B3 = hit(P3)
local A4, B4 = hit(P4)

g:Define_sphere({center=C, radius=R, color="white",opacity=0.6, 
    edgestyle="noline", arrowBstyle="-stealth", show=true})

g:DSpolyline({{O,P1},{O,P2},{O,P3},{O,P4}}, {color="gray",style="solid",arrows=1})
g:DSfacet({A1,A3,A4,A2},{fill="teal",fillopacity=0.8,width=1,style="solid"})
g:DSfacet({B1,B3,B4,B2},{fill="teal",fillopacity=0.8,width=1,style="solid"})
g:Dspherical()

g:Dpolyline3d({{O,4.3*vecI},{O,4.3*vecJ},{O,4.3*vecK}}, "line width=0.6,-stealth")
g:Dpolyline3d({P1,P2,P4,P3}, true)
--g:Dfacet({{A1,A2,B2,B1},{A2,A4,B4,B2},{A4,A3,B3,B4},{A3,A1,B1,B3}}, {mode=5, color="teal", opacity=0.3})
g:Dlabel3d("$x$",4.55*vecI,{pos="SW"},"$y$",4.55*vecJ,{pos="E"},"$z$",4.55*vecK,{pos="N"})
g:Show()
\end{luadraw}
\end{document}
